% ------------------------------------------------------------------------------
% Abstract
% ------------------------------------------------------------------------------

\begin{resumo}[Abstract]

    The Pé-de-Meia Program, established by Brazilian Law No. 14,818/2024, is a savings-based financial incentive for public high school students from families registered in the Unified Registry for Social Programs (\textit{Cadastro Único}), aimed at promoting school retention and completion.
    This paper conducts a nationwide observational assessment of the association between the program's rollout and upper secondary education indicators.
    A balanced panel of 324 observations was built (27 Federative Units $\times$ state and private school systems $\times$ six years between 2019 and 2025), combining payment microdata from the Office of the Comptroller General with school flow rates from the School Census and SAEB/IDEB results from INEP.
    The 2024 financial execution was audited from all nine published monthly payrolls, totalling BRL 5.59 billion in 27.9 million payments, and a correction of 1.79 million duplicated records in the March 2024 payroll was identified.
    Twenty-one difference-in-differences models were estimated with state-by-system and year fixed effects and state-clustered standard errors with a $t$ distribution (26 degrees of freedom).
    In the post-program period (2024--2025), state schools showed, relative to private schools, a differential reduction of 2.07 percentage points in dropout rates and an increase of 4.28 p.p. in pass rates ($p < 0.001$).
    In the 2023 \textit{versus} 2025 window, the state IDEB rose by 0.18 point ($p = 0.0115$), driven mainly by the flow indicator, with no significant gain in standardised proficiency ($p = 0.1830$).
    However, formal tests reject the parallel trends assumption in the pre-program period, which, together with multiple testing and ecological aggregation, precludes causal attribution.
    The results should therefore be read as descriptive associations, and the study provides a reproducible pipeline for future evaluations using individual-level microdata.

    \par\vspace{\baselineskip}

    \noindent\textbf{Keywords}: Pé-de-Meia. Difference-in-differences. School dropout. IDEB. Data Science.
\end{resumo}
